\SourceLeaf{162}
\SixthArticle{30}
From formula 38 it follows that
\[
\frac{\Pi(k,z+\omega)}{\Pi(k,z)}
=\frac{z+1}{z+1+\omega}\cdot\frac{z+2}{z+2+\omega}\cdot\frac{z+3}{z+3+\omega}
\mathinner{.\,.\,.\,.}\frac{z+k}{z+k+\omega}\cdot k^{\omega}
\]
Taking logarithms and expanding them into infinite series gives
\[
\begin{aligned}
{}\text{[60]}\quad&\log\Pi(k,z+\omega)=\log\Pi(k,z)\\
&+\omega\left(\log k-\frac{1}{z+1}-\frac{1}{z+2}-\frac{1}{z+3}-\text{etc.}-\frac{1}{z+k}\right)\\
&+\tfrac12\omega\omega\left(\frac{1}{(z+1)^2}+\frac{1}{(z+2)^2}+\frac{1}{(z+3)^2}+\text{etc.}+\frac{1}{(z+k)^2}\right)\\
&-\tfrac13\omega^3\left(\frac{1}{(z+1)^3}+\frac{1}{(z+2)^3}+\frac{1}{(z+3)^3}+\text{etc.}+\frac{1}{(z+k)^3}\right)\\
&+\text{etc. to infinity}
\end{aligned}
\]
The series multiplied here by $\omega$, which may also be written, if preferred, as
\[
\begin{aligned}
&-\frac{1}{z+1}+\log 2-\frac{1}{z+2}+\log\tfrac32-\frac{1}{z+3}\\[4pt]
&+\log\tfrac43-\frac{1}{z+4}+\log\tfrac54-\text{etc.}
+\log\frac{k}{k-1}-\frac{1}{z+k}
\end{aligned}
\]
consists of a finite number of terms; but as $k$ grows to infinity it will converge to a definite limit, which presents us with a new kind of transcendental function, to be denoted henceforth by $\Psi z$.

Further, denoting the sums of the following series, extended to \textit{infinity},
\[
\begin{array}{l}
\dfrac{1}{(z+1)^2}+\dfrac{1}{(z+2)^2}+\dfrac{1}{(z+3)^2}+\text{etc.}\\[10pt]
\dfrac{1}{(z+1)^3}+\dfrac{1}{(z+2)^3}+\dfrac{1}{(z+3)^3}+\text{etc.}\\[10pt]
\dfrac{1}{(z+1)^4}+\dfrac{1}{(z+2)^4}+\dfrac{1}{(z+3)^4}+\text{etc.}\\[10pt]
\text{etc.}
\end{array}
\]
by $P,Q,R$ etc. respectively (for which it seems less necessary to introduce functional symbols), we shall have
\[
\begin{aligned}
{}\text{[61]}\quad&\log\Pi(z+\omega)=\log\Pi z+\omega\Psi z\\[4pt]
&+\tfrac12\omega\omega P-\tfrac13\omega^3Q+\tfrac14\omega^4R-\text{etc.}
\end{aligned}
\]
Evidently the function $\Psi z$ will be the first derivative of the function $\log\Pi z$, and therefore
\[
{}\text{[62]}\qquad \frac{\mathrm{d}\Pi z}{\mathrm{d}z}=\Pi z\cdot\Psi z
\]
Similarly, $P=\dfrac{\mathrm{d}\Psi z}{\mathrm{d}z}$, $Q=-\dfrac{\mathrm{dd}\Psi z}{2\mathrm{d}z^2}$, $R=+\dfrac{\mathrm{d}^3\Psi z}{2.3\mathrm{d}z^3}$ etc.
\SourceLeaf{163}
\SixthArticle{31}
The function $\Psi z$ is almost as remarkable as the function $\Pi z$; we shall therefore collect its more notable relations here. Differentiating equation 44 gives
\[
\Psi(z+1)=\Psi z+\frac{1}{z+1}\tag*{[63]}
\]
whence
\[
\Psi(z+n)=\Psi z+\frac{1}{z+1}+\frac{1}{z+2}+\frac{1}{z+3}+\text{etc.}+\frac{1}{z+n}\tag*{[64]}
\]
With its help it is possible to proceed from smaller values of $z$ to larger ones, or to work back from larger ones to smaller ones. For larger positive values of $z$, the numerical values of the function can be calculated quite conveniently by the following formulas, obtained by differentiating equations 58 and 59; the same cautions must, however, be kept in mind as those we gave in Article 29 concerning formulas 58 and 59:
\[
\Psi z=\log z+\frac{1}{2z}-\frac{\mathfrak{A}}{2zz}+\frac{\mathfrak{B}}{4z^4}-\frac{\mathfrak{C}}{6z^6}+\text{etc.}\tag*{[65]}
\]
\[
\Psi\left(z-\tfrac12\right)=\log z+\frac{\mathfrak{A}}{2\cdot2zz}-\frac{7\mathfrak{B}}{4\cdot8z^4}+\frac{31\mathfrak{C}}{6\cdot32z^6}-\text{etc.}\tag*{[66]}
\]
Thus, for $z=10$ we calculated
\[
\Psi z=2{,}3517525890\ 6672110764\ 743
\]
from which we work back to
\[
\Psi0=-0{,}5772156649\ 0153286060\ 653\HistoricalNote{18}
\]
For a positive integer value of $z$, in general,
\[
\Psi z=\Psi0+1+\tfrac12+\tfrac13+\text{etc.}+\frac1z\tag*{[67]}
\]
For a negative integer value, however, $\Psi z$ evidently becomes an infinitely large quantity.
\SourceLeaf{164}
\SixthArticle{32}
Formula 55 gives us $\log\Pi(-z)+\log\Pi(z-1)=\log\pi-\log\sin z\pi$, from which differentiation gives
\[
\Psi(-z)-\Psi(z-1)=\pi\operatorname{cotang}z\pi\tag*{[68]}
\]
And since the definition of the function $\Psi$ gives, in general,
\[
\Psi x-\Psi y=-\frac1{x+1}+\frac1{y+1}-\frac1{x+2}+\frac1{y+2}-\frac1{x+3}+\text{etc.}\tag*{[69]}
\]
the known series follows:
\[
\pi\operatorname{cotang}z\pi=\frac1z-\frac1{1-z}+\frac1{1+z}-\frac1{2-z}+\frac1{2+z}-\frac1{3-z}+\text{etc.}
\]
Similarly, differentiating formula 57 gives
\[
\begin{aligned}
&\Psi z+\Psi\left(z-\frac1n\right)+\Psi\left(z-\frac2n\right)+\text{etc.}+\Psi\left(z-\frac{n-1}{n}\right)\\[4pt]
&=n\Psi nz-n\log n
\end{aligned}
\tag*{[70]}\]
and hence, on setting $z=0$,
\[
\begin{aligned}
&\Psi\left(-\frac1n\right)+\Psi\left(-\frac2n\right)+\Psi\left(-\frac3n\right)+\text{etc.}+\Psi\left(-\frac{n-1}{n}\right)\\[4pt]
&=(n-1)\Psi0-n\log n
\end{aligned}
\tag*{[71]}\]
Thus, for example,
\[
\begin{aligned}
&\Psi\left(-\tfrac12\right)=\Psi0-2\log2=-1{,}9635100260\ 2142347944\ 099,\\[4pt]
&\quad\text{and hence further}
\end{aligned}
\]
\[
\Psi\tfrac12=+0{,}0364899739\ 7857652055\ 901.
\]
